\subsection{学习率调度}

现代 LLM 训练几乎都使用学习率调度器（learning rate scheduler），而不是固定学习率。最常见的策略是 \textbf{warmup + cosine decay}：

\begin{enumerate}
    \item \textbf{Warmup 阶段}：从接近 0 的学习率线性增长到峰值学习率，通常持续总步数的 1\%--5\%。
    \item \textbf{Cosine decay 阶段}：学习率按余弦函数从峰值衰减到一个很小的最终值（通常为峰值的 1/10 或更小）。
\end{enumerate}

$$
\eta_t = \begin{cases}
\eta_{\max} \cdot \frac{t}{T_{\text{warmup}}} & \text{if } t < T_{\text{warmup}} \\[6pt]
\eta_{\min} + \frac{1}{2}(\eta_{\max} - \eta_{\min})\left(1 + \cos\left(\frac{t - T_{\text{warmup}}}{T_{\text{total}} - T_{\text{warmup}}} \cdot \pi\right)\right) & \text{otherwise}
\end{cases}
$$

\begin{knowledgebox}{为什么需要 warmup}
训练初期，Adam 的二阶矩 $v_t$ 尚未充分预热（bias correction 后的估计仍不稳定），此时使用大学习率容易导致梯度爆炸或训练发散。Warmup 让模型在参数空间中先”小步试探”，等优化器状态稳定后再加大步长。

经验上，warmup 步数通常设为总步数的 1\%--2\%，或者 2000 步左右。
\end{knowledgebox}

